\section{工业4.0下人的角色}

\subsection{工业4.0下人的角色}

在一个到处都是自动化、机器人、传感器的工业4.0世界里，人到底还有没有位置？答案其实是肯定的，我们将从概念、身体负荷到具体案例,一步步说明为什么人依然重要,以及技术应该怎么帮人。

在工业4.0之前,90年代很流行一个叫CIM(计算机集成制造)的理念,核心思路是尽量把人从生产线上“计算机化”掉。理想状态就是一座\emph{不需要人的工厂}，\emph{计算机占据主导地位}，人只是生产中次要的角色。工业4.0则反过来了。\emph{人被重新放回中心位置}，角色变成了问题解决者、\emph{经验和知识的载体}、做决策的人。\emph{生产这一侧则交给自动化}、包容性设计、人体增强、各种辅助系统，以及能察言观色、根据情境自动调整的交互界面去承担。

要理解为什么需要辅助系统,得先看清人在车间里实际承受的身体负荷。不同的工作类型会集中损耗身体的不同部位：负重搬运主要伤腰椎、肩膀、髋关节和膝盖；头顶作业主要伤颈椎、肩膀、胸椎和腰椎；精密装配则容易造成肩膀、颈部、手腕的劳损，还需要持续的注意力。

工业4.0里人的角色发生了根本转变，\emph{从CIM时代“能不用人就不用人”,变成了“人是核心,技术要配合人”；再用身体负荷的分析，论证辅助系统存在的必要性；然后用一堆真实案例证明这些技术已经在各行各业落地；最后指出目前这个领域还缺一套系统的分类方法。}

\subsection{人机分工}

四次工业革命，每一次都会重新定义人需要具备什么能力。我们把这些能力归纳成六个能力领域：
\begin{itemize}
    \item 身体负荷承受能力(physische Beanspruchung)
    \item 情绪负荷承受能力(emotionale Beanspruchung)
    \item 心理/脑力负荷承受能力(mentale Beanspruchung)
    \item 专业能力(Fachkompetenz)
    \item 社会能力(Sozialkompetenz)
    \item 个性能力(Persönlichkeitskompetenz)
\end{itemize}
也就是说，工业4.0对人的要求已经不只是“体力好、手艺熟”,而是身体、心理、专业、社交、个性多方面综合的能力。

既然要谈谁该做什么,就得先弄清楚人和机器分别擅长什么。人的优势集中在认知能力(Kognition)、感知运动协调能力(Sensomotorik)、学习能力(Lernfähigkeit)。简单来说人擅长的就是灵活应变、理解情境、不断学习。

技术系统的优势集中在:精度(Genauigkeit)、速度(Geschwindigkeit)、可靠性(Zuverlässigkeit)、抗负荷能力(Belastbarkeit,包括脑力和体力两方面)、耐力和抗疲劳/抗分心能力 (Ausdauer/~Ermüdungsanfälligkeit/Konzentrationsfähigkeit)，所以机器擅长的是稳定、精确、不知疲倦地重复。

既然人和机器的强项在很多地方是相反但互补的，那就应该通过合理搭配，用技术系统、工具和功能去补上人的短板，同时保留人在认知和灵活应变上的优势，而不是简单地用机器把人整个替换掉。这就是人机智能混合(intelligente Hybridisierung von Mensch und Technik)。

\paragraph{任务分配的阶段(Stufen der Aufgabenverteilung)}

如图 \ref{fig:stufenaufgabenverteilung} 所示，从Stufe 0(完全人工)到Stufe 4(完全自动化)。我们把一项工作拆成四个维度,看每个维度在不同阶段里是由人还是由机器主导:
\begin{itemize}
    \item 过程主导权(Prozesshoheit):谁来决定整个流程怎么走
    \item 监控(Überwachung):谁来盯着任务执行得对不对
    \item 感知(Sensorik):谁来获取环境和任务状态的信息
    \item 施力(Kraftaufbringung):谁来实际出力完成动作
\end{itemize}

\begin{figure}[htbp]
  \centering
  \includegraphics[width=0.8\textwidth]{stufenaufgabenverteilung.png}
  \caption{任务分配的阶段}
  \label{fig:stufenaufgabenverteilung}
\end{figure}

有意思的是,这四个维度并不是同步从人切换到机器的,而是有先后顺序:最先交给机器的是施力(比如用电动工具代替纯手力),然后是感知(比如用传感器代替人眼观察),接着是监控(比如系统自动报警代替人工巡检),最后才轮到过程主导权(也就是让机器自己决定怎么做),这也是自动化程度最高、最难实现的一步。两端分别对应完全人工作业(vollmanuelle Tätigkeit)和完全自动化(Vollautomatisierung),中间则是各种程度不一的人机协作组合。

\subsection{技术支持系统}

人和人是不一样的,所以支持系统不能是一刀切的,必须足够灵活。

造成个体差异的因素被分成了两类。一类是人本身的特征:年龄、性别、既往经验、认知与注意力水平、动作模式与行为习惯。另一类是可以量化的身体参数:人体测量数据、力量、柔韧性、心肺功能/耐力、既往损伤史。这些因素叠加起来,意味着同一套支持系统,对不同的人可能效果完全不同。这也是后面模块化设计思路的起点。

各个活动场景需要的支持方式不一样，有的需要力的转向,有的需要力的转向加放大,有的需要单纯的力量放大,还有的需要稳定与质量保证。所以支持系统不应该是一个笼统的概念,而是要针对具体场景设计具体功能。

我们用三个维度来判断一个人机协作属于哪一种形式：
\begin{itemize}
    \item \emph{时空关系}:人和机器是分开工作、共处一地但独立、分布式协作,还是完全融合?
    \item \emph{耦合形式}:是辅助还是更紧密的耦合?
    \item \emph{控制权归属}:由人还是由技术主导,而且这种主导是暂时性的还是永久性的?
\end{itemize}
这三个维度组合起来,划分出三条典型路径,分别对应不同紧密程度的人机协作关系。这也是理解后面各种具体支持系统"到底算哪一类"的一个坐标系。

一个技术系统要被称为支持系统,必须同时满足四个条件:
\begin{enumerate}
    \item 它能帮助人完成任务,而不是完全或部分地替代人
    \item 执行的主导权仍然掌握在人手里(由操作者给出目标值,而不是系统强制规定)
    \item 人是这个系统的操作者
    \item 系统本身不会对操作者或第三方造成危险
\end{enumerate}

除了减轻当下的身体负荷，引入支持系统还有三层更长远的理由：
\begin{itemize}
    \item \emph{补偿机能衰退} (Kompensation von Funktionseinbußen)：人的体力、柔韧性、视力、反应速度都会随年龄增长而下降，既往损伤也会留下永久性的能力缺口。支持系统可以把这部分损失的能力补回来，让因年龄或身体条件受限的人依然能胜任原本的工作岗位，而不是被迫调岗或离开生产线。
    \item \emph{更长久地保持健康} (längerer Erhalt der Gesundheit)：负重搬运、头顶作业、精密装配这类工作的损伤是长期累积的，等到肌肉骨骼疾病显现时往往已经不可逆。支持系统的作用是提前把负荷峰值压下去，把作用在腰椎、肩关节上的力转移到设备上，从而延长人的职业寿命，降低企业的病假成本和工伤风险。
    \item \emph{技术接受度与可接受性} (Technikakzeptanz und -akzeptabilität)：一套系统再先进，员工不愿意用也等于零。接受度指的是使用者主观上愿不愿意用——设备是否舒适、是否拖慢节奏、是否让人觉得被监视；可接受性则是更外部的判断，即这套系统在伦理、法律和社会层面是否应当被接受——数据采集的边界在哪里、会不会造成新的能力依赖。只有两者同时成立，支持系统才可能真正落地。
\end{itemize}

\paragraph{支持系统的模块化设计思路}

支持系统必须因人而异、因任务而异,\emph{为了实现这种灵活性}，我们把支持系统拆成几类可以自由组合的模块。

首先是基础参数模块,涵盖人体\emph{测量数据、负荷、能力、技能、复杂度、持续时间,用于描述这个人和这项任务的具体需求}。比如在头顶高度用重工具装配时,需要弯曲/伸展/旋转功能加抬起/放下/保持功能,再配合膝部支撑等具体模块组合出一套解决方案。

其次是人机接口模块,按感知渠道分成触觉、视觉、听觉三类。

然后是传感器模块,分成两类接口:物理接口(比如EMG肌电传感器、一维/三维力传感器、力敏薄膜传感器、压力传感器)和认知接口(比如摄像系统、眼动追踪、脑电EEG传感器、语音识别)。

最后是控制与调节模块,同样分物理和认知两类支持:物理支持包括力控制(按位置调节或由使用者主动调节)和意图识别(无意识/自动化控制,结合EMG传感器和动作识别);认知支持包括防错机制(Poka-Yoke)、工作现场监控(比如用摄像系统监控作业、偏差提醒)、轨迹/抓取路径给定(从CAD数据或系统学习中获得),以及反馈功能(声音、视觉、触觉多种形式提示偏差)。

这一整套模块化思路的意义在于与其为每个人、每个任务单独设计一套系统,不如把常用的功能拆成标准化的模块,按需组合。

\paragraph{支持也是一把双刃剑}

支持系统并不是越多越好, 而是有代价的。

好处：
\begin{itemize}
    \item \emph{}身体减负 (körperliche Entlastung)：外骨骼、力平衡装置把原本压在腰椎和肩关节上的负荷转移到设备上, 肌肉骨骼疾病明显减少。
    \item \emph{}质量提升 (Steigerung der Qualität)：防错机制、轨迹给定和实时反馈让作业结果更稳定,出错次数下降。
    \item \emph{}生产率提升 (Steigerung der Produktivität)：力量放大和辅助定位缩短了单件工时, 人还能获得超能力式的力量加成。
    \item \emph{}动作追踪(Bewegungsverfolgung)：支持系统本身就装有力传感器、EMG 传感器或摄像系统, 这些数据顺带就能用于工效学评估、作业分析和长期健康监测, 不需要额外部署一套采集系统。
\end{itemize}

代价：
\begin{itemize}
    \item \emph{接受度问题} (Akzeptanzprobleme)：设备穿戴不舒服、拖慢了原本的作业节奏、或者让员工觉得自己被监视, 都会导致系统被闲置。
    \item \emph{系统设计错误}(falsche Systemauslegung)：支持力给得不对、支撑点选得不对, 不仅达不到减负效果,反而可能把负荷转移到别的部位, 制造出新的劳损。
    \item \emph{岗位轮换被取消}(Aussetzen der Job-Rotation)：一旦某个工位配了专用支持设备, 企业往往倾向于让固定的人长期留在这个工位上, 原本靠轮岗实现的负荷分散就没有了, 单一部位的累积负荷反而上升。
    \item \emph{能力退化}：长期依赖支持设备可能导致肌肉退化、身体能力下降; 认知层面同理, 辅助给得太多会让人逐渐丧失原本具备的判断力。所以支持的度要拿捏得当, 这件事从长期看非常重要。支持不足解决不了问题,
\end{itemize}

